% !TeX program = xelatex
% ============================================================
% project_proposal/project_proposal.tex — the course project description
% and application, 1st year of the HSE MSc programme «System and Software
% Engineering» (English edition of the mentor seminar form, rev. 2.1).
%
% Preliminary edition (PRP, due with the topic application by 15 November):
% the starred sections are required. Extended edition (EPP, checkpoint CP1
% in January): adds the abstract, terms, abbreviations, work plan, resources
% and risks. Switch the edition in the project settings («Proposal edition»)
% — the \ifhseExtendedProposal blocks below appear without editing the text.
%
% Replace the yellow \hseFill{…} with your text. Your data, the supervisor
% and the titles come from the project settings.
% ============================================================
\documentclass[11pt,a4paper]{article}
% !TeX program = xelatex
% ============================================================
% project_proposal/preamble.tex — преамбула «Описания проекта и заявки
% на выполнение» курсового проекта магистратуры СПИ.
%
% Повторяет вёрстку шаблона семинара наставника (ред. 2.1): A4, поля
% 20/15/20/20 мм, синие заголовки разделов, таблица «Сводка». Документ
% подаётся в PDF вместе с заявлением о теме (SmartLMS, до 15 ноября).
% Язык разделов — язык работы (project.lang).
% ============================================================
\newcommand{\hseDefaultMono}{Consolas}
\input{../common/fonts}
((* if project.lang == 'en' *))
\usepackage[russian,english]{babel}
((* else *))
\usepackage[english,russian]{babel}
((* endif *))

\usepackage[a4paper,left=20mm,right=15mm,top=20mm,bottom=20mm]{geometry}
\usepackage{setspace}
\setstretch{1.15}
\usepackage{indentfirst}
\setlength{\parindent}{1.25cm}
\setlength{\parskip}{4pt}

\usepackage{microtype}
\usepackage{csquotes}
\usepackage{amsmath}
\usepackage{amssymb}
\usepackage{etoolbox}
\usepackage{graphicx}
\usepackage{xcolor}
\usepackage{array}
\usepackage{tabularx}
\usepackage{booktabs}
\usepackage{longtable}
\usepackage{enumitem}
\setlist{nosep, leftmargin=1.25cm}
\setlist[itemize,1]{label=--}
\usepackage{xurl}
\urlstyle{same}
\usepackage{hyperref}
\hypersetup{unicode=true,colorlinks=true,linkcolor=black,urlcolor=black,citecolor=black}
\usepackage{cite}
\usepackage{lastpage}
\usepackage{fancyhdr}

% ── Цвета и заголовки разделов (как в шаблоне семинара) ─────
\definecolor{hseprpnote}{RGB}{230,184,183}
\usepackage{titlesec}
\titleformat{\section}{\large\bfseries}{\thesection.}{0.4em}{}
\titlespacing*{\section}{0pt}{14pt}{6pt}
\titleformat{\subsection}{\normalsize\bfseries}{}{0em}{}
\titlespacing*{\subsection}{0pt}{10pt}{4pt}
% Обязательный раздел (поле со звёздочкой в шаблоне ОП).
\newcommand{\hseReq}{\textsuperscript{*}}

% ── Колонтитул: номер страницы и редакция формы ─────────────
\pagestyle{fancy}
\fancyhf{}
\renewcommand{\headrulewidth}{0pt}
\fancyfoot[C]{\thepage}
\fancyfoot[R]{\scriptsize ((* if project.lang == 'en' *))Form rev. 2.1 of 1.11.2023((* else *))Ред. 2.1 от 1.11.2023((* endif *))}

% Номера в списке источников с точкой.
\makeatletter
\renewcommand\@biblabel[1]{#1.}
\makeatother

\input{../common/meta}

\begin{document}

\begin{center}
{\Large\bfseries\color{hseprpblue} Course Project}\\[2pt]
{\color{hseprpblue}(1st year of the MSc programme «System and Software Engineering»)}\\[10pt]
{\Large\bfseries\color{hseprpblue} Project Description and Application}\\[2pt]
{\color{hseprpblue}\ifhseExtendedProposal(extended)\else(preliminary)\fi}
\end{center}

\subsection*{Summary\hseReq}
\noindent
\renewcommand{\arraystretch}{1.35}
\begin{tabularx}{\textwidth}{|>{\raggedright\arraybackslash\bfseries}p{0.3\textwidth}|>{\raggedright\arraybackslash}X|}
\hline
1. Title in Russian & \foreignlanguage{russian}{\hseTitleRu} \\
\hline
2. Title in English & \hseTitleEn \\
\hline
3. Language of the work & \hseWorkLangEn \\
\hline
4. Student (author) & \hseAuthorName, group \hseGroup \\
\hline
5. Supervisor: name, academic title, degree and position & \hseSupervisorFull
\ifhseHasCoSupervisor\newline Consultant: \hseCoSupervisorName\ifdefempty{\hseCoSupervisorDegree}{}{, \hseCoSupervisorDegree}\ifdefempty{\hseCoSupervisorTitle}{}{, \hseCoSupervisorTitle}\fi \\
\hline
6. Supervisor's signature and date &
\vspace{1.6cm}
% Put the supervisor's signature in the Signatures tool (PNG) or by hand;
% with a detached e-signature (GnuPG) replace the line below with
% «Signed with a detached electronic signature».
\rule{0.8cm}{0.4pt}.\rule{0.8cm}{0.4pt}.\the\year \\
\hline
\end{tabularx}

\ifhseExtendedProposal
\subsection*{Abstract}
% One or two paragraphs (500–1500 characters): what the work is about —
% problem, approach, expected result. Not a summary of results.
\hseFill{what the work is about: problem, approach, expected result}
\fi

\setcounter{section}{6}
\section{Subject area, relevance and motivation\hseReq}
% Subject area: context and boundaries. Relevance: why it matters now and
% to whom. Motivation: the research gap the work closes and why you.
\hseFill{the subject area and the boundaries of the work}

\hseFill{why the topic matters now: who benefits from the solution}

\hseFill{motivation: which gap in existing solutions the work closes}

\subsection*{Related work and existing solutions\hseReq}
\hseFill{existing approaches and solutions with citations and their limitations}~\cite{denscombe2019}.

\subsection*{Problem statement\hseReq}
\hseFill{the specific problem the work solves}

\noindent\textbf{Research question:} \hseFill{the question the work answers}.

\noindent\textbf{Object of study:} \hseFill{the broad area, process or phenomenon the work addresses}.

\noindent\textbf{Subject of study:} \hseFill{the specific part, property or aspect of the object under study}.

\section{Keywords}
% Keywords are separated by semicolons (;).
\hseFill{keyword; keyword; keyword}

\section{Aim of the work\hseReq}
% One aim: action verb (develop, determine, evaluate) + object + by what
% means. SMART.
\hseFill{develop … for … based on …}

\section{Tasks}
% 3–7 steps from analysis to implementation and evaluation; they form the
% basis of the work plan — do not merge them into coarse items.
\begin{enumerate}[label=\arabic*)]
  \item \hseFill{analyse existing solutions and approaches};
  \item \hseFill{define the requirements};
  \item \hseFill{design the architecture (method, model)};
  \item \hseFill{implement a prototype};
  \item \hseFill{test and evaluate the results}.
\end{enumerate}

\section{Expected results\hseReq}
% Results are MEASURABLE: it must be easy to check whether they are
% achieved; give the verification method (metric, threshold) if needed.
\begin{enumerate}[label=\arabic*)]
  \item \hseFill{result} --- \hseFill{how to verify: metric and threshold};
  \item \hseFill{result} --- \hseFill{how to verify}.
\end{enumerate}

\section{Methods}
\hseFill{expected methods}

\section{Tools}
\hseFill{expected tools}

\subsection*{Additional sections (optional)}
\ifhseExtendedProposal
\medskip\noindent\textbf{Terms and definitions.}\par
\begin{description}[leftmargin=0pt, font=\bfseries]
  \item[\hseFill{term}] --- \hseFill{definition}.
\end{description}

\medskip\noindent\textbf{Abbreviations.}\par
\begin{description}[leftmargin=0pt, font=\bfseries]
  \item[\hseFill{ABBR}] --- \hseFill{expansion}.
\end{description}

\medskip\noindent\textbf{Work plan.}\par
\noindent
\begin{tabularx}{\textwidth}{|p{0.45\textwidth}|p{0.15\textwidth}|X|}
\hline
\textbf{Stage} & \textbf{Due} & \textbf{Deliverable} \\
\hline
\hseFill{task 1} & \hseFill{month} & \hseFill{artefact} \\
\hline
\hseFill{task 2} & \hseFill{month} & \hseFill{artefact} \\
\hline
\end{tabularx}

\medskip\noindent\textbf{Resources and risks.}\par
\hseFill{resources, constraints and main risks with responses}
\else
\hseOptional{Any further information needed to present the idea of the course project, for example appendices.}
\fi

\renewcommand{\refname}{References\hseReq}
\begin{thebibliography}{99}
  \bibitem{denscombe2019}
  M.~Denscombe, \emph{Research Proposals: A Practical Guide}, 2nd~ed. London: Open University Press, 2019.
\end{thebibliography}

\end{document}
